% !PRINTED 299
\weberpage{299}
\begin{center}
\textbf{\S\ 81. \gesperrt{Nichtnormale Teiler von $\mathfrak{L}_0$.}}
\end{center}

Wir fragen nun, indem wir den Fall $p = 3$ beiseite lassen, nach den nichtnormalen Teilern der Gruppe $\mathfrak{L}_0$, deren Index kleiner als $p+1$ ist; von der Existenz solcher Teiler hängt die Möglichkeit der Bildung von Resolventen der Transformationsgleichung ab, \gesperrt{deren Grad niedriger ist als $p+1$.}

Zunächst läßt sich zeigen, daß der Index eines Teilers $\mathfrak{K}$ von $\mathfrak{L}_0$ niemals kleiner als $p$ sein kann.

Es sei
\weq{1}{\mathfrak{K} = S_0,\ S_1,\ \ldots,\ S_{\nu-1}}
irgend ein Teiler von $\mathfrak{L}_0$ vom Grade $\nu$, und, wenn es möglich ist,
\weq{2}{T = \begin{pmatrix} \alpha, & \beta \\ \gamma, & \delta \end{pmatrix}}
eine Substitution $p$ten Grades, die in $\mathfrak{L}_0$, aber nicht in $\mathfrak{K}$ vorkommt. Es kann dann auch, da $p$ Primzahl ist, keine niedrigere Potenz von $T$ als die $p$te in $\mathfrak{K}$ vorkommen. Die Nebengruppen
\weq{3}{\mathfrak{K},\ T\mathfrak{K},\ T^2\mathfrak{K},\ \ldots,\ T^{p-1}\mathfrak{K}}
enthalten lauter \gesperrt{voneinander verschiedene Elemente} $T^i S_k$, und demnach ist $\nu p$ höchstens gleich dem Grade der Gruppe $\mathfrak{L}_0$, d.\,h.:
\[
\nu p \leqq p\,\frac{p^2-1}{2}.
\]
Also ist der Index von $\mathfrak{K}$, d.\,h. der Quotient
\[
p\,\frac{p^2-1}{2} : \nu
\]
\gesperrt{gleich $p$ oder größer als $p$.}

Ein Teiler $\mathfrak{K}$ von $\mathfrak{L}_0$, dessen Index kleiner als $p$ ist, muß daher sämtliche Substitutionen $p$ten Grades enthalten, also auch (§ 79, I.) alle Substitutionen $T$, in denen
\weq{4}{\alpha + \delta \equiv 2 \pmod{p}}
ist. Demnach enthält eine solche Gruppe $\mathfrak{K}$ zunächst die Substitution
\weq{5}{A = \begin{pmatrix} 1, & 0 \\ 1, & 1 \end{pmatrix}}
%CONT>>
